\section{Proof Sketch}

First form the seed-averaged terminal responses over all complete public
response rules, and choose a basis of their finite-dimensional span before
selecting any learning instance.  Lemma~\ref{lem:step-001-complete-rules} and
Propositions~\ref{prop:step-001-mean-response-space} and
\ref{prop:step-001-fixed-coordinates} make this construction precise and
separate mean-response rank from the generally much larger span of individual
seed-specific predictors.

Next choose finitely many point evaluations that coordinatize \(V_A\).  They
put every raw mean response in a common bounded set, so the closed convex body
\(K_A\) is compact in one fixed topology and every point evaluation is
continuous.  For each distribution-target pair, the exact-center complete
rule is tolerance-valid on every reached seed path.  The universal guarantee
and the binary loss identity then produce a member of \(K_A\) with signed
correlation at least \(\rho\).

For a nonempty finite set \(S\), play the bilinear payoff
\(L_h(f,p)=\sum_{x\in S}p(x)h(x)f(x)\) on the fixed body \(K_A\) and the
simplex \(\Delta(S)\).  Sion's theorem changes the order of optimization,
while the exact-center witness may vary with \(p\) inside the same fixed body.
This yields one \(f_{h,S}\) with margin at least \(\rho\) at every point of
\(S\).  The resulting pointwise constraint sets are closed subsets of the
same compact \(K_A\) and have the finite-intersection property, including the
empty subfamily.  Compactness therefore gives a single \(f_h\) satisfying all
points of an arbitrary domain.

Finally, exact basis coordinates turn \(f_h\) into \(w_h\) with zero score
residual.  Strict positivity makes \(r_A\) admissible in the definition of
dimension complexity, and the primitive rank certificate gives the displayed
polynomial upper bound.  The appendix proves each of these statements in full,
including the empty, zero-rank, zero-query, noiseless, and parameter-boundary
cases.
